\section{Each Element in Its Setting}
\label{sec:each-element}

The texts of this Mass come from two books with separate histories. The
antiphons and orations belong to the Missal's Twenty-sixth Week in Ordinary
Time and are prayed with the readings of all three years; the readings, psalm
and acclamation belong to the Lectionary's Year~A. Each element has its own
literary setting and textual facts, and the witnesses who expounded it read it
in that setting.

\subsection{The confession from the furnace}

The Entrance Antiphon, \latin{Omnia, quae fecisti nobis, Domine}, is drawn from
the prayer that Azarias, one of the three young men thrown into
Nabuchodonosor's furnace for refusing to adore the golden statue, prays ``in
the midst of the fire'' (Dan 3:25). The prayer belongs to the part of Daniel 3
that the Greek Bible and the Vulgate carry and the Hebrew does not. Jerome
marks the place in his commentary: \latin{Hucusque Hebraei legunt}, the Hebrews read thus far; what follows, to the end of the song of the
three young men, \latin{in Hebraico non habentur}, they do not have in the
Hebrew (\work{In Danielem}, on 3:23, PL 25, 509). The prayer opens by blessing
God as just in all he has done, his ``ways right, and all thy judgments
true'' (3:27), confesses that the people ``have sinned, and committed
iniquity'' and ``have not hearkened to thy commandments'' (3:29--30), pleads
for God's name and covenant (3:34--36), and, with no holocaust or sacrifice
left, asks that ``in a contrite heart and humble spirit'' the people be
accepted as though they were holocausts of rams and bullocks (3:39--40).

The antiphon takes five of these verses, in the order 31, 29, 30, 43 and 42,
and makes one short confession and one plea out of them. Its words are not the
Vulgate's, as the account of the appointed text shows, and Ildefonso Schuster
observed of the same chant that the verses are Daniel's ``but the actual words
of the text are not quoted'' (\work{The Sacramentary}, vol.~3, p.~174). The
speaker is striking. Azarias and his companions are in the furnace because
they refused to sin, yet they say ``we have sinned''. Jerome saw the point at
once: \latin{Et certe tres pueri non peccaverant}, and certainly the three
young men had not sinned; they speak \latin{ex persona populi}, in the person
of the people, as Paul speaks of evil he does not will in Romans~7. On the
opening verse he draws the lesson for any hearer: \latin{Quando diversis
premimur angustiis, ex toto cordis hoc loquamur affectu, et quidquid nobis
acciderit, iuste nos sustinere fateamur}, when we are pressed by distresses,
let us say this with the whole heart, and confess that whatever has befallen
us we bear justly (\work{In Danielem}, PL 25, 509).

\subsection{The Collect: power shown in sparing}

The Collect, \latin{Deus, qui omnipotentiam tuam}, closes the Introductory
Rites after the Gloria. Its address makes a claim about God that would sound
strange if it were not so familiar: that he shows his almighty power above all
in sparing and in mercy. The prayer is old. Its Latin stands word for word in
the Gregorian Sacramentary in the form edited by H.~A. Wilson as the
\work{Hadrianum} (1915, p.~172), and, with one possible variant, in the Old
Gelasian Sacramentary. Thomas Aquinas cited it as what \latin{de Deo dicitur}, is said of God, and
set it among the objections in his article on God's omnipotence
(\work{Summa theologiae} I, q.~25, a.~3). His reply gives three reasons why
the prayer is right. \latin{Dei omnipotentia ostenditur maxime in parcendo et
miserando, quia per hoc ostenditur Deum habere summam potestatem, quod libere
peccata dimittit}: God's omnipotence is shown most in sparing and showing
mercy, because it shows that he has the highest power that he freely forgives
sins. Or else because, \latin{parcendo hominibus et miserando}, by sparing men
and showing them mercy, \latin{perducit eos ad participationem infiniti boni,
qui est ultimus effectus divinae virtutis}, he leads them to a share in the
infinite good, which is the last effect of the divine power. Or because the
effect of divine mercy is the foundation of all God's works, since nothing is
owed to anyone except on the ground of what God first gave unowed (ad~3). In
another place Aquinas adds that mercy is, taken in itself, the greatest of the
virtues, because it pours itself out on others and relieves their defects, and
that belongs to the superior; \latin{unde et misereri ponitur proprium Deo, et
in hoc maxime dicitur eius omnipotentia manifestari}, and so to have mercy is
set down as proper to God, and his omnipotence is said to be shown most in it
(II-II, q.~30, a.~4).

Schuster, the Benedictine Cardinal Archbishop of Milan, called
the Collect ``of truly classic beauty'' and took its thought further: ``the
rehabilitation of one who has gone astray in his sins requires, so to speak, a
condescension and powerful energy on the part of God, such as the creation of
the world itself never once called for'', for the abyss between God and evil
is deeper than that between him and nothing, and God ``bridges over this abyss
when in his infinite mercy he descends into it to draw forth the sinner''
(\work{The Sacramentary}, vol.~3, pp.~121--122). The petition that follows asks
for grace upon those who run, \latin{currentes}, toward God's promises.

\subsection{Ezekiel 18: the exiles' proverb and the Lord's way}

The first reading is taken from a chapter that answers a proverb. Ezekiel, a
priest among the captives ``by the river Chobar'' (Ezek 1:1--3), confronts a
saying repeated ``in the land of Israel'': ``The fathers have eaten sour
grapes, and the teeth of the children are set on edge'' (18:2). The disaster
that had fallen on the people was, on this view, the bill for their ancestors'
sins, and nothing they did could change it.
The Lord's answer is that ``all souls are mine'' and ``the soul that sinneth,
the same shall die'' (18:4, 20). The chapter works this out in a series of
cases: a just father, a violent son, a grandson who sees his father's sins and
does not imitate them (18:5--18); then the wicked man who does penance and
lives, of whom God says, ``I will not remember all his iniquities that he hath
done'' (18:21--22), and the just man who turns to iniquity, of whom ``all his
justices which he hath done, shall not be remembered'' (18:24). Between the
two cases stands the question on which the chapter turns: ``Is it my will that
a sinner should die, saith the Lord God, and not that he should be converted
from his ways, and live?'' (18:23).

The appointed verses begin with the people's objection, as the Lord reports
it. \latin{Non est aequa
via Domini}, the Vulgate has it: the Lord's way is not level, not fair. The
Lord turns the charge back on them, ``are not rather your ways perverse?'', and
restates the two cases (18:26--28). The same complaint returns at v.~29, just
beyond the appointment, and the chapter ends with the call that gives the whole
its purpose: ``Be converted, and do penance \ldots\ make to yourselves a new
heart, and a new spirit: and why will you die, O house of Israel? \ldots\ return
ye and live'' (18:30--32). The chapter carries no date of its own; in the book
it falls between the oracle dated to the sixth year of the captivity (8:1) and
the one dated to the seventh (20:1). The \work{Ordo} prefixes the prophetic
formula \latin{Haec dicit Dominus} to v.~25, which in the text begins by
reporting the people's complaint: \latin{Et dixistis}, ``And you have said''.

\subsection{Psalm 25 (24): the prayer of one who asks the way}

The psalm is an alphabetic prayer in the Hebrew, and the Vulgate and the
Douay--Rheims head it ``Unto the end, a psalm for David''. Its speaker is beset by enemies and burdened by old sins,
and he asks for three things: to be shown the Lord's ways, to be remembered
according to mercy, and to be taught. The appointed strophes gather the first
two requests and the confidence behind them: ``Shew, O Lord, thy ways to me,
and teach me thy paths'' (v.~4); ``Remember, O Lord, thy bowels of compassion
\ldots\ The sins of my youth and my ignorances do not remember'' (vv.~6--7);
``The Lord is sweet and righteous: therefore he will give a law to sinners in
the way. He will guide the mild in judgment: he will teach the meek his ways''
(vv.~8--9). The word ``way'' recurs through the psalm, at vv.~4, 8, 9, 10 and
12, and the verse after the appointment sums it up: ``All the ways of the Lord
are mercy and truth'' (v.~10). The response is the first half of v.~6 alone.
Two textual details bear on what is sung. The Vulgate's v.~4 opens with a
clause on the confusion of those who act unjustly, which the Lectionary's
numbering places in v.~3; and the last clause of v.~5, ``on thee have I waited
all the day long'', may stand outside the strophe as the United States
Lectionary prints it.

\subsection{Philippians 2: an appeal and a hymn}

The second reading begins with a ``therefore'' (\latin{ergo}) that the
\work{Ordo}'s opening words omit. It joins the appeal to what precedes: Paul has
just asked the Philippians to ``stand fast in one spirit, with one mind
labouring together for the faith of the gospel'', undaunted by adversaries,
since to them ``it is given for Christ, not only to believe in him, but also to
suffer for him'' (1:27--30). The appeal of 2:1--4 is for one mind and one
charity, for nothing done ``through contention'' or ``vain glory'', and for
each to ``esteem others better than themselves'' and look to ``those that are
other men's''. Verse~5 names its ground: ``let this mind be in you, which was
also in Christ Jesus''. The longer form continues with the hymn of vv.~6--11:
Christ, ``being in the form of God, thought it not robbery to be equal with
God'', but emptied himself, took the form of a servant, humbled himself and became
obedient unto death, ``even to the death of the cross''; therefore God exalted
him and gave him the name above all names, at which every knee shall bow and
every tongue confess. The chapter goes on at once to draw the consequence,
``with fear and trembling work out your salvation'', ``without murmurings'', in
the midst of ``a crooked and perverse generation'' (2:12--15).

The Lectionary offers the two forms for a pastoral reason: the introduction to
the \work{Ordo} leaves the choice to the hearers' capacity to listen with profit
to a longer or a shorter reading (\work{Praenotanda} 80). The choice is
material. The shorter form keeps the appeal and names the mind of Christ, which
is exactly what the \work{Ordo}'s title for the reading, \latin{Sentite in vobis
quod et in Christo Iesu}, singles out; the account of what that mind did
belongs to the longer form alone.

\subsection{The Alleluia verse: the sheep who hear}

The acclamation sings a verse of the Good Shepherd discourse as John places it
at the feast of the Dedication in Jerusalem, in winter, in Solomon's porch
(Jn 10:22--23). Those who gathered round Jesus asked him to say plainly whether
he was the Christ. He answered, ``I speak to you, and you believe not'', and,
``you do not believe, because you are not of my sheep'' (10:25--26). The
appointed verse follows: ``My sheep hear my voice. And I know them: and they
follow me.'' The next verses promise eternal life and say that none can snatch
the sheep from his hand or his Father's (10:28--30). The Lectionary adds a
speaker formula to the verse, and the \work{Ordo}'s Latin has \latin{dicit
Dominus}. Verse~26, which the
acclamation does not sing, has a history of its own in the Fathers.

\subsection{Matthew 21: a question about John}

The Gospel stands between two other answers to the same challenge. When
Jesus was teaching in the temple, ``the chief priests and ancients of the
people'' asked him by what authority he acted (Mt 21:23). He asked them in
turn whether John's baptism was from heaven or from men, and they would not
say: from heaven, and they would be asked why they had not believed him; from
men, and they feared the people, who held John for a prophet (21:24--27). The
parable of the two sons follows at once, and its application returns to John:
``John came to you in the way of justice: and you did not believe him'' (21:32).
The parable of the wicked husbandmen comes next (21:33--46), and at its end the
chief priests and Pharisees ``knew that he spoke of them'' (21:45). The
parable's first hearers are therefore the men who would not answer the
question about John, and it is spoken to draw their answer from their own
mouths.

The two sons are drawn with a few strokes. The first says ``I will not'', and
``afterwards, being moved with repentance, he went''; the second says ``I go,
Sir'', with a title of respect, and does not go. When Jesus applies the story
he uses the words of the first son for his hearers, and in the negative: they
``did not even afterwards repent'' (21:32). The Latin makes the echo plain,
\latin{postea \ldots\ poenitentia motus} of the son and \latin{nec
poenitentiam habuistis postea} of the chief priests. The \work{Ordo} sets both
halves of the saying over the reading, the son's repentance and the publicans'
precedence, and prints the latter in the present tense, \latin{praecedunt},
where the Clementine Vulgate has the future \latin{praecedent}.

The order of the sons and the answer of v.~31 are not the same in every
witness to the text. The Lectionary's text, as the bishops' conference prints
it for the day, has the father go first to the son who refuses and afterwards
goes and then to the one who assents and does not go, and has the hearers
answer by naming the first; the \work{Nova Vulgata} (\latin{Primus}), the
Clementine and the Douay--Rheims agree. So, according to the apparatus of the
Society of Biblical Literature's Greek New Testament, do the twenty-eighth
Nestle--Aland edition, the Robinson--Pierpont text and the SBL edition itself.
The same apparatus reports that Westcott and Hort print the sons the other way
round, the first promising and not going and the second refusing and
repenting, and with Tregelles give the answer ``the latter''. The Latin Fathers
knew a divided text as well. Jerome quotes the hearers' answer as
\latin{Novissimus}, the last, and then says, \latin{Sciendum est in veris
exemplaribus non haberi novissimum, sed primum}, that the true copies read not
``the last'' but ``the first'', \latin{ut proprio iudicio condemnentur}, so that
they are condemned by their own judgment. He explains either reading: if
\latin{novissimus} is read, the Jews understood the truth \latin{sed
tergiversari, et nolle dicere quod sentiunt}, but evaded it and would not say
what they thought, as they would not say that John's baptism was from heaven
(\work{Commentarii in Matthaeum} III, PL 26, 156). His editor Vallarsi notes
that Rabanus Maurus read \latin{in veteribus}, ``in the old copies'', for
\latin{in veris}. Hilary of Poitiers worked with a text that has the
Lectionary's order of the sons but in which the Pharisees answer by naming the
younger, and his Maurist editor notes that this answer was the current reading
in Hilary's and Jerome's time (\work{Commentarius in Matthaeum} 21, PL 9,
1039--1041). Hilary's identification of the two sons rests on that answer.

\subsection{The Prayer over the Offerings}

The Prayer over the Offerings, \latin{Concede nobis, misericors Deus}, asks the
merciful God that the assembly's offering be acceptable and that through it
\latin{fons omnis benedictionis}, the fountain of every blessing, be opened to
them. It ends with the short conclusion. Its plea for acceptance has a
counterpart in the prayer from which the Entrance Antiphon comes, where
Azarias, with no holocaust or sacrifice left to offer, asks that ``in a
contrite heart and humble spirit'' the people themselves be accepted, ``so let
our sacrifice be made in thy sight this day, that it may please thee'' (Dan
3:38--40).

\subsection{The two Communion antiphons}

The first antiphon, \latin{Memento verbi tui servo tuo, Domine}, is marked
\latin{Cf.} and adapts two verses of the long alphabetic psalm on God's law, at
the strophe of the letter Zain. In the psalm the servant asks God to be mindful
of the word ``in which thou hast given me hope'', and says that this comforted
him ``in my humiliation: because thy word hath enlivened me'' (Ps 118:49--50);
the next verse sets him beside the proud who ``did iniquitously altogether''
(118:51). The antiphon ends before the last clause of v.~50. Its
\latin{Memento} is not the Clementine's \latin{Memor esto} but an older Latin
psalter reading: Hilary heads his exposition of the strophe \latin{Memento
uerbi tui seruo tuo, in quo mihi spem dedisti}, in the antiphon's word order,
and Augustine's lemma has \latin{Memento} too (\work{Tractatus super Psalmos},
CSEL 22, p.~418; \work{Enarrationes in Psalmos} 118, s.~15.1).

The second antiphon, printed under \latin{Vel}, is 1~John 3:16 as the
Clementine has it: ``In this we have known the charity of God, because he hath
laid down his life for us: and we ought to lay down our lives for the
brethren.'' In the letter the verse stands between the warning against Cain,
who ``killed his brother'', and the test of charity that follows it: ``He that
hath the substance of this world and shall see his brother in need and shall
shut up his bowels from him: how doth the charity of God abide in him? My
little children, let us not love in word nor in tongue, but in deed and in
truth'' (1~Jn 3:12, 17--18). The Missal's rubric for Ordinary Time prefers,
between two Communion antiphons, the one that agrees with the Gospel of the
Mass; neither antiphon of this pair is drawn from a Gospel, and the books do
not say whether either agrees with this one.

\subsection{The Prayer after Communion}

The Prayer after Communion, \latin{Sit nobis, Domine, reparatio mentis et
corporis}, asks that the heavenly mystery be the restoring of mind and body. Its
first clause is old: it stands in the Hadrianum, in the Old Gelasian and in the
Verona collection of Mass prayers, where it ends differently, \latin{ut cuius
exsequimur actionem sentiamus effectum}, that we may feel the effect of the
action we perform. The present ending, which joins the proclaiming of the
Lord's death and the sharing of his sufferings to a share in his glory as
fellow heirs, stands in none of the three. Its words \latin{coheredes} and
\latin{compatimur} stand together in Romans 8:17, ``heirs indeed of God and
joint heirs with Christ: yet so, if we suffer with him, that we may be also
glorified with him'', and \latin{mortem \ldots\ annuntiando} recalls 1~Corinthians
11:26, ``you shall shew the death of the Lord, until he come''. The prayer
closes with \latin{Qui vivit et regnat}, because its last clause speaks of the
Son.

\subsection{Three families of words}

Three families of words run through these texts, in the Latin of the Clementine
Vulgate and in the Douay--Rheims. The first is the Lord's way: \latin{via
Domini} in Ezekiel's complaint (18:25, 29), \latin{vias tuas} and \latin{vias
suas} in the psalm (24:4, 9), and, in the prayer from which the Entrance
Antiphon comes, \latin{viae tuae rectae} (Dan 3:27). The second is remembering
and not remembering: God will not remember the penitent's iniquities nor the
lapsed man's justices (Ezek 18:22, 24); the psalmist asks God to remember his
compassion and not the sins of youth (Ps 24:6--7); the servant asks him to
remember his word (Ps 118:49). The third is humility: \latin{in humilitate}
in Paul's appeal and \latin{humiliavit semetipsum} in the hymn (Phil 2:3, 8),
the \latin{mites} whom the Lord teaches (Ps 24:9), the servant's comfort
\latin{in humilitate mea} (Ps 118:50), and the \latin{spiritu humilitatis} of
Azarias's prayer (Dan 3:39). One witness joins two of these texts himself.
Expounding the Communion psalm's \latin{Memento}, Augustine asks whether God can
forget, and answers from Scripture's own way of speaking, citing among others
\latin{Omnes \ldots\ iniquitates eius obliviscar} from Ezekiel 18:22
(\work{Enarrationes in Psalmos} 118, s.~15.1).
